\documentclass[12pt,a4paper]{article}

% --- Настройка языка и шрифтов для LuaLaTeX ---
\usepackage{fontspec}
\defaultfontfeatures{Ligatures=TeX,Scale=MatchLowercase}

\setmainfont{Alice-Regular}[
    Path=font/,
    Extension=.ttf,
    BoldFont = Alice-Regular,
    ItalicFont = Alice-Regular,
    BoldItalicFont = Alice-Regular,
    BoldFeatures = {FakeBold=1.6},
    ItalicFeatures = {FakeSlant=0.2},
    BoldItalicFeatures = {FakeBold=1.6,FakeSlant=0.2}
]

\newfontfamily\cyrillicfont{Alice-Regular}[
    Path=font/,
    Extension=.ttf,
    BoldFont = Alice-Regular,
    ItalicFont = Alice-Regular,
    BoldItalicFont = Alice-Regular,
    BoldFeatures = {FakeBold=1.6},
    ItalicFeatures = {FakeSlant=0.2},
    BoldItalicFeatures = {FakeBold=1.6,FakeSlant=0.2}
]

\usepackage{polyglossia}
\setmainlanguage{russian}
\setsansfont{TeX Gyre Heros}
\setmonofont{TeX Gyre Cursor}

\usepackage{graphicx}        					% Графика
\usepackage{amsmath, amsfonts, amssymb, amsthm}	% Математика
\usepackage{mathtools}
\usepackage{chemfig}
\usepackage[version=4,arrows=pgf-filled,
            textfontname=sffamily,
            mathfontname=mathsf]{mhchem}
\usepackage[protrusion=false]{microtype}       % Микротипографика
\usepackage{braket} 						% Braket нотация для квантовой механики
\usepackage{unicode-math} 					% Современный пакет для математики в LuaLaTeX
\usepackage{longtable}       				% Таблицы с разрывом страниц
\usepackage{tikz}           				% Рисование
\usepackage[europeanresistors]{circuitikz}	% Рисование электрических схем в европейском стиле
\usepackage{listings}						% Для отображения кода
\usepackage{tabularx}       				% Для таблиц во всю ширину
\usepackage{array}		  					% Для расширенных возможностей таблиц
\usepackage{cancel}                % Перечёркивать
\usepackage{float}          				% Для точного позиционирования фигур
\usepackage{pgfplots}                          % Для построения графиков
\usepackage{caption}                           % Для настройки подписей к рисункам и таблицам
\usepackage{xfp}                               % Для вычислений в LaTeX
\usepackage{luacode}                           % Для выполнения Lua-кода в документе    
\pgfplotsset{compat=1.18}
\usepackage[paper=letterpaper,
            top=20mm,
            bottom=20mm,
            left=10mm,
            right=10mm,
            footskip=0.4in 
]{geometry}
\usepackage{colortbl}

\usetikzlibrary{arrows.meta}

% --- Колонтитулы ---
\usepackage{fancyhdr}
\pagestyle{fancy}
\fancyhf{}

% --- Теоремы и окружения ---
\theoremstyle{plain}
\newtheorem{theorem}{Theorem}[section]
\newtheorem{lemma}[theorem]{Lemma}
\newtheorem{proposition}[theorem]{Proposition}
\newtheorem{corollary}[theorem]{Corollary}

\theoremstyle{definition}
\newtheorem{definition}[theorem]{Definition}

\theoremstyle{remark}
\newtheorem{remark}[theorem]{Remark}

\numberwithin{equation}{section}

\renewcommand{\headrulewidth}{0.4pt}
\renewcommand{\footrulewidth}{0.4pt}
\renewcommand\lstlistingname{Исходный код}

\lstdefinestyle{main}{
	backgroundcolor=\color{gray!8},   
	commentstyle=\color{green!50!black},
	keywordstyle=\color{blue!80!black},
	numberstyle=\small\color{darkgray},
	stringstyle=\rmfamily\color{orange!90!black},
	basicstyle=\ttfamily\small,
	breakatwhitespace=false,         
	breaklines=true,                   
	keepspaces=true,                 
	numbers=left,                    
	numbersep=8pt,                  
	showspaces=false,                
	showstringspaces=false,
	showtabs=false,                  
	tabsize=4
}	

\lstset{style=main}

\newcolumntype{C}{>{\centering\arraybackslash}X}
\renewcommand{\arraystretch}{1.25}

\begin{luacode*}
include = require("include")
\end{luacode*}

\begin{document}

\begin{center}
    \Large Лабораторная работа \number 3
\end{center}

\begin{center}
    \large Реактивное ионное травление тонких диэлектрических пленок
\end{center}

\textbf{Цель работы}: освоение технологии реактивного ионного травления (РИТ) тонких диэлектрических пленок; изучение влияния параметров процесса РИТ на селективность травления двуокиси кремния и нитрида кремния; получение практических навыков работы на установке МВУ ТМ РИТ $100$.

\section{Теоретические сведения}

\subsection{Физическая сущность РИТ и возникновение автосмещения}

Реактивное ионное травление занимает промежуточное положение между изотропным плазмохимическим травлением (ПХТ), где съем материала идет нейтральными химически активными радикалами при низких энергиях частиц ($< 10$~эВ), и чисто физическим ионным распылением, в котором тяжелые ионы инертного газа выбивают атомы решетки кинетическим ударом с энергиями до нескольких килоэлектронвольт. В РИТ травление осуществляется совокупным действием химически активных радикалов ($\mathrm{F}^\bullet$, $^\bullet\mathrm{CF}_3$, $\mathrm{Cl}^\bullet$) и направленного потока ускоренных реактивных ионов ($\mathrm{CF}_3^+$, $\mathrm{CF}_2^+$, $\mathrm{SF}_5^+$, $\mathrm{Cl}_2^+$) с энергиями в диапазоне $100\text{--}500$~эВ.

Процесс реализуется в диодных ВЧ-реакторах (частота генератора $13{,}56$~МГц) с емкостной связью. Пластина размещается на изолированном от земли катоде, подключенном через разделительный конденсатор, а стенки рабочей камеры служат заземленным анодом. Из-за колоссальной разницы в массах подвижность электронов в плазме на порядки превосходит подвижность ионов ($\mu_e \gg \mu_i$). В первые четвертьпериоды ВЧ-колебаний электроны мгновенно устремляются к катоду, заряжая его отрицательно относительно квазинейтральной плазмы. Чтобы суммарный ток через изолированный электрод за период ВЧ-поля был равен нулю, на катоде формируется постоянный отрицательный потенциал самосмещения $U_\text{см}$:
\begin{equation*}
  U_\text{см} \approx \frac{V_\text{ВЧ}}{2}, \qquad \frac{V_c}{V_a} \approx \left(\frac{A_a}{A_c}\right)^4,
\end{equation*}
где $V_c$ и $V_a$ — падения потенциала в приэлектродных темных пространствах катода и анода, а $A_a$ и $A_c$ — эффективные площади заземленных стенок камеры и катодного столика. Поскольку стенки камеры много больше столика подложки ($A_a \gg A_c$), почти вся разность потенциалов падает в тонком слое пространственного заряда непосредственно над катодом.

«Если вы просто зажжете разряд без отрицательного смещения на столике, вы получите чисто изотропную химическую ванну, только в газовой фазе. Вся суть РИТ держится на том, что разделительный конденсатор превращает ВЧ-поле в постоянное запирающее смещение на сотни вольт, которое вытягивает ионы из плазмы и бросает их строго в лоб пластине».

\subsection{Синергизм процессов и эксперимент Кобурна-Винтерса}

Главное достоинство РИТ — неаддитивность, или синергический эффект: скорость комбинированного травления многократно превышает арифметическую сумму скоростей чисто химического травления и чисто физического распыления:
\begin{equation*}
  V_\text{РИТ} \gg V_\text{хим} + V_\text{физ}.
\end{equation*}

Классическим экспериментальным доказательством синергизма стал опыт Гарольда Винтерса и Джона Кобурна (1979 г.). При подаче на поверхность кремния изолированного молекулярного пучка дифторида ксенона $\mathrm{XeF}_2$ скорость химического травления составляла порядка $5$~\AA/мин. При воздействии только пучка ионов аргона $\mathrm{Ar}^+$ с энергией $450$~эВ скорость физического распыления не превышала $2$~\AA/мин. Однако при их одновременном включении суммарная скорость травления кремния скачкообразно возрастала до $50\text{--}100$~\AA/мин — более чем на порядок.

Физический механизм синергизма складывается из взаимного усиления двух процессов:
\begin{enumerate}
  \item Хемосорбция активных радикалов галогенов ослабляет ковалентные связи поверхностных атомов кремния с подложкой, снижая поверхностный потенциальный барьер и многократно увеличивая коэффициент физического распыления $K_\text{р}$ поверхностных комплексов.
  \item Удар ускоренного иона передает импульс атомам решетки, генерируя локальные радиационные дефекты, разрывая связи $\mathrm{Si}\text{--}\mathrm{Si}$ и $\mathrm{Si}\text{--}\mathrm{O}$ на глубину нескольких монослоев. В этой разупорядоченной зоне химическая реакция с радикалами протекает практически без барьера активации, а локальный тепловой всплеск и импульсный удар вызывают взрывную десорбцию промежуточных фторидов, очищая поверхность для следующего цикла реакции.
\end{enumerate}

По характеру взаимодействия с радикалами реакции разделяют на ионно-ускоренные (химический процесс между материалом и ХАЧ идет и спонтанно, но ионы форсируют его скорость, как при травлении $\mathrm{Si}$ фтором) и ионно-стимулированные, когда при комнатной температуре без ионной активации реакция вообще термодинамически заблокирована, что фундаментально характерно для диэлектриков.

\subsection{Химия травления диэлектриков ($\mathrm{SiO}_2$, $\mathrm{Si}_3\mathrm{N}_4$) и механизм селективности}

В отличие от чистого кремния, энергия ковалентной связи в оксиде кремния составляет $\sim 800$~кДж/моль ($\mathrm{Si}\text{--}\mathrm{O}$), в то время как связь $\mathrm{Si}\text{--}\mathrm{Si}$ требует всего $\sim 327$~кДж/моль, а связь $\mathrm{Si}\text{--}\mathrm{N}$ в нитриде кремния занимает промежуточное положение ($\sim 439$~кДж/моль). Из-за высокой прочности связи термические радикалы фтора $\mathrm{F}^\bullet$ при комнатной температуре не способны спонтанно разрушать решетку $\mathrm{SiO}_2$: скорость спонтанного травления оксида фтором практически равна нулю.

Травление $\mathrm{SiO}_2$ возможно только при одновременном присутствии ионной стимуляции и фторуглеродных радикалов $^\bullet\mathrm{CF}_x$ ($x = 1, 2, 3$). Падающие ионы разбивают связи $\mathrm{Si}\text{--}\mathrm{O}$, а хемосорбированные радикалы $\mathrm{CF}_2$ забирают кислород решетки, образуя летучие монооксид и диоксид углерода:
\begin{equation*}
  \mathrm{SiO}_2 + 2\mathrm{CF}_2 \xrightarrow{\text{ионы}} \mathrm{SiF}_4\uparrow + 2\mathrm{CO}\uparrow \quad (\text{или } \mathrm{CO}_2\uparrow).
\end{equation*}

Травление нитрида кремния $\mathrm{Si}_3\mathrm{N}_4$ протекает с образованием летучего молекулярного азота или смешанных цианистых соединений:
\begin{equation*}
  \mathrm{Si}_3\mathrm{N}_4 + 12\mathrm{F}^\bullet \xrightarrow{\text{ионы}} 3\mathrm{SiF}_4\uparrow + 2\mathrm{N}_2\uparrow.
\end{equation*}

Селективность травления $S = V_{\text{тр}_1} / V_{\text{тр}_2}$ регулируют управлением балансом полимеризации и травления через соотношение концентраций фтора и углерода ($\mathrm{F}/\mathrm{C}$) в газовой фазе. При использовании фреонов, обедненных фтором ($\mathrm{CHF}_3$, $\mathrm{C}_4\mathrm{F}_8$, $\mathrm{C}_2\mathrm{F}_6$) или при добавлении в плазму водорода $\mathrm{H}_2$, водород необратимо связывает свободные радикалы фтора в прочную молекулу $\mathrm{HF}$:
\begin{equation*}
  \mathrm{CF}_4 \to \mathrm{CF}_3^\bullet + \mathrm{F}^\bullet, \qquad \mathrm{F}^\bullet + \mathrm{H}^\bullet \to \mathrm{HF}\uparrow.
\end{equation*}

Падение концентрации фтора сдвигает равновесие в сторону осаждения защитной тефлоноподобной полимерной пленки $(\mathrm{CF}_2)_n$. На открытой поверхности кремния эта пленка непрерывно нарастает и наглухо блокирует дальнейший доступ травящих агентов к подложке: травление $\mathrm{Si}$ полностью останавливается. На поверхности же $\mathrm{SiO}_2$ кислород, освобождаемый при разрыве связей решетки, окисляет осаждающийся углерод в летучие газы $\mathrm{CO}$ и $\mathrm{CO}_2$, самостоятельно очищая поверхность оксида. В результате $\mathrm{SiO}_2$ продолжает активно травиться при нулевой скорости травления кремния, что обеспечивает высокую селективность $S(\mathrm{SiO}_2/\mathrm{Si}) \approx 15\text{--}30$.

«Кислород сидит прямо в узлах кристаллической решетки оксида и сам выжигает углеродный нагар, превращая его в летучий угарный газ. А кремнию взять кислород неоткуда: он моментально затягивается полимерной коркой и засыпает. Отсюда и берется избирательность оксида к подложке».

Однако при повышении энергии ионов свыше $1$~кэВ полимерный слой с кремния начинает сбиваться прямым физическим распылением, и селективность резко падает. Для травления нитрида кремния $\mathrm{Si}_3\mathrm{N}_4$ используют фреоно-водородные смеси с добавками кислорода: кислород связывает избыток углерода и освобождает фтор, повышая скорость травления $\mathrm{Si}_3\mathrm{N}_4$ относительно кремния и оксида.

\subsection{Магнетронный разряд и технологические особенности МВУ ТМ РИТ 100}

В классических диодных реакторах для поддержания разряда при давлениях $1\text{--}10$~Па требуется подавать высокую ВЧ-мощность. Это вызывает рост напряжения автосмещения до $400\text{--}600$~В: ускоренные ионы вызывают глубокие радиационные дефекты подзатворного диэлектрика, внедряются в решетку кремния и разрушают фоторезистивную маску за счет перегрева.

Радикальным решением стало внедрение магнетронного реактивного ионного травления, реализованного в отечественной установке МВУ ТМ РИТ $100$ (модульная вакуумная установка, технологический модуль РИТ для подложек диаметром до $100$~мм). 

Под плоским водоохлаждаемым катодом размещается кольцевая магнитная система на постоянных самарий-кобальтовых или неодимовых магнитах, создающая над поверхностью катода арочное магнитное поле с индукцией $B \approx 0{,}03\text{--}0{,}05$~Тл ($300\text{--}500$~Гс). Силовые линии поля перпендикулярны вектору электрического поля $\vec{E}$ приэлектродного слоя. В скрещенных полях $\vec{E} \times \vec{B}$ вторичные электроны захватываются в замкнутую магнитную ловушку и совершают циклоидальный дрейф вдоль поверхности катода. Длина траектории электрона до ухода на анод возрастает в сотни раз, многократно увеличивая вероятность ионизирующих столкновений с молекулами газа.

Это дает установке МВУ ТМ РИТ $100$ три решающих преимущества:
\begin{enumerate}
  \item Устойчивое горение плотной плазмы при рекордно низких рабочих давлениях ($0{,}05\text{--}0{,}5$~Па), где длина свободного пробега ионов превышает ширину прикатодной темной зоны ($\overline{\lambda_i} > d_\text{тп}$). Ионы преодолевают слой без столкновений, падая на пластину строго ортогонально.
  \item Высокая скорость травления ($0{,}5\text{--}1{,}2$~мкм/мин) при умеренной энергии ионов ($E_i \le 150\text{--}250$~эВ). Отсутствует радиационный пробой тонких подзатворных оксидов и исключается термическая деструкция маски.
  \item Эффективное водяное охлаждение катодного столика и откачка турбомолекулярным насосом с прецизионной системой напуска газов ($\mathrm{CF}_4$, $\mathrm{CHF}_3$, $\mathrm{SF}_6$, $\mathrm{O}_2$, $\mathrm{Ar}$), обеспечивающие воспроизводимость травления по пластине с разбросом не более $\pm 3\%$.
\end{enumerate}

«Магнитное поле не ускоряет ионы — оно закручивает электроны в спираль и удерживает их у катода. Электроны не бьются о стенки камеры, а бесконечно ионизируют газ на месте. В итоге мы получаем густую плазму при давлении в десять раз ниже обычного, а ионы летят прямо и аккуратно, не разбивая кристалл».

\subsection{Механизмы анизотропии и паразитные эффекты субмикронного РИТ}

Анизотропия процесса количественно характеризуется отношением скоростей травления:
\begin{equation*}
  A = \frac{v_\perp}{v_\parallel} = \frac{d}{b}, \qquad A_n = 1 - \frac{v_\parallel}{v_\perp} = 1 - \frac{1}{A},
\end{equation*}
где $d$ — глубина вытравленной канавки, $b$ — величина бокового подтравливания под маску, $A$ — коэффициент анизотропии, а $A_n$ — нормированный показатель анизотропии ($A_n = 0$ для чисто изотропного съема, $A_n \to 1$ для строго вертикальных стенок). В установке МВУ ТМ РИТ $100$ достигаются значения $A \approx 20\text{--}100$ ($A_n \ge 0{,}95\text{--}0{,}99$).

Вертикальность профиля в РИТ гарантируется совместным действием двух факторов:
\begin{enumerate}
  \item \textbf{Направленная ионная активация дна}: вектор движения ускоренных ионов строго ортогонален поверхности подложки. Ионы стимулируют химическую реакцию только на горизонтальном дне канавки, в то время как вертикальные стенки практически не испытывают ионных ударов.
  \item \textbf{Ингибиторная пассивация боковых граней}: полимеризующиеся фрагменты газовой смеси ($(\mathrm{CF}_2)_n$) равномерно оседают на всей поверхности рельефа. На дне канавки пленка непрерывно стравливается и распыляется налетающими ионами, а на боковых стенках она сохраняется и служит барьером, защищающим диэлектрик от травления блуждающими радикалами фтора.
\end{enumerate}

При уменьшении проектных норм ниже $0{,}5$~мкм начинают проявляться специфические паразитные эффекты РИТ:
\begin{enumerate}
  \item \textbf{Зависимость скорости от аспектного отношения (ARDE / RIE-lag)}: канавки с высоким отношением глубины к ширине ($d / w > 3$) травятся существенно медленнее широких окон. Это вызвано кнудсеновским диффузионным сопротивлением узких каналов для нейтральных радикалов и геометрическим затенением ионного потока стенками маски.
  \item \textbf{Микромаскирование («плазменная трава»)}: атомы нелетучих конструкционных материалов катода или камеры, а также фрагменты запекшегося резиста оседают на дно канавки в виде субмикронных кластеров. Они выступают локальными микромасками, под которыми вырастают плотные вертикальные иглы кремния или оксида.
  \item \textbf{Дифференциальный заряд диэлектриков (notching)}: электроны плазмы обладают изотропным распределением скоростей и оседают на верхней кромке диэлектрической маски, заряжая ее отрицательно. Ионы же летят вертикально и заряжают изолирующее дно положительно. Возникающее локальное электрическое поле искривляет траектории низкоэнергетических ионов к основанию боковой стенки, формируя клиновидную выемку (подсечку) на границе раздела диэлектрика с кремниевой подложкой.
\end{enumerate}

\subsection{Критерии перетравливания и технология самоформирования спейсеров}

При травлении ступенчатого топологического рельефа интегральной схемы номинального времени травления $t_d = d / v_\text{п}$ недостаточно: из-за технологического разброса толщины пленки по пластине ($\delta_d \approx \pm 5\%$) и флуктуаций скорости травления ($\delta_v \approx \pm 5\%$) на ступеньках рельефа остается недотравленный материал. Для гарантированного вскрытия окон вводят технологическое перетравливание:
\begin{equation*}
  t_\text{п} = t_d (1 + f),
\end{equation*}
где $f = 0{,}1\text{--}0{,}5$ ($10\text{--}50\%$) — коэффициент перетравливания. 

Необходимость перетравливания накладывает жесткие аналитические ограничения на селективность процесса:
\begin{equation*}
  S(\text{слой}/\text{маска}) \ge 4 \frac{d_\text{п}}{d_\text{м}} f (1 + \delta_d + \delta_v),
\end{equation*}
\begin{equation*}
  S(\text{слой}/\text{подслой}) \ge \frac{d_\text{п}}{\Delta d_\text{под}} \left[-1 + (1 + f)(1 + \delta_d + \delta_v)\right],
\end{equation*}
где $d_\text{п}$ и $d_\text{м}$ — толщины травящейся пленки и фоторезиста, а $\Delta d_\text{под}$ — предельно допустимая потеря толщины нижележащего изолирующего подслоя (например, тонкого подзатворного оксида).

Показатель анизотропии травления пленки должен строго удовлетворять условию сохранения проектных топологических норм:
\begin{equation*}
  A \ge \frac{f d_\text{п} (1 + \delta_d + \delta_v)}{\Delta L_\text{пв}},
\end{equation*}
где $\Delta L_\text{пв}$ — допустимый технологический уход линейного размера элемента по верхней кромке.

Высокая вертикальная анизотропия РИТ лежит в основе технологии самоформирования субмикронных элементов (технологии спейсеров), позволяющей обходить дифракционный предел проекционной фотолитографии:
\begin{enumerate}
  \item На исходной подложке фотолитографией и анизотропным травлением создается базовая вертикальная ступенька оксида кремния $\mathrm{SiO}_2$.
  \item Методом химического осаждения из газовой фазы (CVD) конформно наносится тонкая пленка поликристаллического кремния или нитрида толщиной $d_\text{конф}$. На вертикальной боковой грани ступеньки эффективная толщина слоя вдоль нормали к плоскости подложки удваивается.
  \item Проводится безмасочное анизотропное РИТ (процесс etch-back) строго на глубину осажденной толщины $d_\text{конф}$. На плоских планарных участках пленка полностью стравливается до подложки, а вдоль вертикального уступа ступеньки остается узкая самоформированная полоска материала — пристеночный спейсер.
\end{enumerate}

Ширина полученного спейсера задается исключительно толщиной осажденной CVD-пленки $d_\text{конф}$, контроль которой на атомном уровне легко обеспечивается временем осаждения. В прямом методе спейсер из поликремния используется непосредственно как субмикронный затвор транзистора шириной $30\text{--}50$~нм. В косвенном методе спейсер служит жесткой маской для травления нижележащих проводящих слоев, позволяя кратно увеличивать плотность интеграции СБИС без перехода на дорогостоящее экстремальное ультрафиолетовое (EUV) оборудование.

«Ширину затвора в $30$ нанометров никакой рядовой оптический степпер на пластине не нарисует — дифракция размоет свет в пятно. Но если вырастить конформный слой толщиной ровно в $30$ нанометров на торце ступеньки, а потом сдуть все лишнее направленным ионным пучком, на стенке останется полоска ровно в толщину слоя. Это чистая топологическая геометрия, обошедшая дифракционный предел оптики».

\end{document}
